\cotefile{60}
\scanpage{36}\mspage{6 - 31 -}
\hfill n$^{\circ}$ 392

que $\operatorname{div}_{B}(Z)$ n'est pas l'image canonique d'un diviseur de $D(A)$. (Raisonner par l'absurde en supposant que $e(\underline{P}/\underline{p})=1$ pour $\underline{P}=BZ$, $\underline{p}=P\cap A$ ; il y aurait alors dans $A$ une uniformisante de $B_{\underline{P}}$, nécessairement de la forme $a(X,Y,Z)Z$ avec $b=a(X,Y,0)\ne 0$ ; en déduire que l'on aurait $\Delta b/b=-XY$ et obtenir une contradiction en calculant $\Delta(-XY)$).

24) \qquad a) Soient $L$ un corps de caractéristique $2$, $\Delta$ une dérivation de $L$, $K$ le sous-corps de $L$ noyau de $\Delta$ ; on suppose que $[L:K]=2$. Montrer que l'on a $\Delta^2=a\Delta$ avec $a\in K$, et que pour qu'un élément $t\in L$ soit de la forme $\Delta x/x$, il faut et il suffit que $\Delta t=at+t^2$.

b) Soient $B$ un anneau local factoriel de caractéristique $2$, $\underline{m}$ son idéal maximal, $L$ son corps des fractions, $\Delta$ une dérivation de $L$ ; on suppose que le sous-corps $K$ de $L$, noyau de $\Delta$, est tel que $[L:K]=2$ ; en outre, on suppose qu'il existe deux éléments $x$, $y$ de $\underline{m}$ tels que $\Delta x$ et $\Delta y$ engendrent l'idéal $\underline{q}$ de $B$ engendré par $\Delta(B)$. Montrer alors que si $t=\Delta z/z$ appartient à $\underline{q}$, il existe un élément inversible $u$ de $B$ tel que $t=\Delta u/u$, (écrire $t=r\Delta x+s\Delta y$ avec $r$, $s$ dans $B$ et utiliser a)).

25) Soient $k$ un corps de caractéristique $2$, $B$ l'anneau de séries formelles $k[[X,Y]]$, $L$ son corps des fractions, $\Delta$ la $k$-dérivation de $L$ définie par $\Delta(X)=Y^{2j}$, $\Delta(Y)=X^{2i}$, ($i$, $j$ entiers $\geq 0$) ; le sous-anneau $A$ de $B$ formé des $x\in B$ tels que $\Delta x=0$ est l'anneau des séries formelles de $k[[X,Y,Z]]$ où on substitue $X^2$ à $X$, $Y^2$ à $Y$ et $X^{2i+1}+Y^{2j+1}$ à $Z$.

a) Montrer que le groupe $C(A)$ contient un espace vectoriel sur $k$ de dimension $N(i,j)$, égale au nombre des couples d'entiers $(a,b)$ tels que $0\leq a<i$, $0\leq b<j$ et $(2j+1)a+(2i+1)b\geq 2ij$. (Utilisant l'exerc. 23 b) et l'exerc. 24 a), noter que les éléments de $L$ sont les séries formelles $F\in B$ telles que $\Delta F=F^2$ ; en attribuant à $X$ le poids $2j+1$ et à $Y$ le poids $2i+1$, décomposer $F$ en somme infinie de polynômes isobares ; si $L_q$ est le sous-groupe de $L$ formé des $F\in L$ dont les composantes isobares sont de poids $\geq q$, $L/L'$ est isomorphe à la somme directe des groupes $C_q/C_{q+1}$, où $C_q=L_q/(L'\cap L_q)$ ; calculer ces groupes pour $q\geq 4ij$, en utilisant l'exerc. 24 b)).

b) Montrer que l'idéal $AX^2$ de $A$ est premier ; si $A'=A[X^{-2}]$, en déduire que $C(A')$ et $C(A)$ sont isomorphes, (n$^{\circ}$\qquad, prop.\qquad). Montrer que $A'$ est un anneau de Dedekind, (considérer l'anneau $B'=B[X^{-2}]$, qui est entier sur $A'$ et est principal, et utiliser le chap. V, \S~2, n$^{\circ}$ 4, th. 3). En déduire un exemple d'anneau de Dedekind dont le groupe des classes d'idéaux est infini.
% archive sheet number: 35
\scanpage{37}
On peut songer à le rendre additif, en le définissant d'abord sur les $E_{\unsure{\eta|y}}$ simples (voire \unsure{abst.|abs\textsuperscript{t}.|abs.} simples, quitte à se permettre des extensions finies sur $\unsure{\mathbb{Q}_{\ell}|\Phi_{\ell}}$) et en prolongeant par linéarité, ce qui revient à définir des facteurs locaux d'un $L_{E_\eta}(\unsure{t|T})$ par $L_{E_\eta(x)}(t)$, en posant $E_\eta(x)=\sum E_\eta^{i*}(x)$, les $E_\eta^i$ étant les facteurs simples \unsure{dans|de} une suite de composition de $E$. Ceci n'est pas encore satisfaisant, car les $E_\eta(x)$ ainsi construits ne sont pas de caractère local en $(E,x)$, --- i.e. ne commutent pas à l'extension \unsure{étale} \deleted{\unsure{sur|en} $X$} (à ext. résiduelle \inserted{en $x$} \deleted{\unsure{en}} triviale) \unsure{sur|en} $X$.
% archive sheet number: 36
\scanpage{38}\mspage{6 - 27 -}
\hfill n$^{\circ}$ 392

12) Démontrer la prop. 5 du n$^{\circ}$ 8 en utilisant le chap. III, \S~2, n$^{\circ}$ 8, cor. 3 du th. 1.

13) Etendre le cor. de la prop. 7 du n$^{\circ}$ 8 au cas où l'anneau local séparé et complet $A$ n'est pas intègre, (utiliser \uline{Alg}., chap. VIII, \S~6, exerc. 6 b)).

14) Soient $A$ un anneau local complet noethérien, dont le corps résiduel est de caractéristique $p>0$. Dans l'anneau de séries formelles $A[[T]]$, on considère les éléments $\omega_n=(1-T)^{p^n}$ et $\gamma_n=1-\omega_n$ pour tout entier $n>0$. Montrer que $\gamma_n$ est, au signe près, un polynôme distingué, (n$^{\circ}$ 8) ; en déduire que $A_n=A[[T]]/(\gamma_n)$ s'identifie à l'algèbre sur $A$ du groupe $G_n=\underline{Z}/p^n\underline{Z}$. Montrer que l'intersection des idéaux principaux $(\gamma_n)$ est réduite à $0$ ; en déduire que $A[[T]]$ s'identifie à la limite projective $\varprojlim A_n$.

15) Soient $K$ un corps complet pour une valuation discrète $v$, $A$ l'anneau de la valuation, $k$ son corps résiduel, $P$ un polynôme de $A[X_1,\ldots,X_n]$ de degré total $d$ ayant la propriété suivante : il existe une extension algébrique $K'$ de $K$ telle que dans $K'[X_1,\ldots,X_n]$, $P$ soit produit de polynômes de degré total $1$. On suppose en outre qu'il existe deux polynômes $Q$, $R$ de $A[X_1,\ldots,X_n]$ tels que $Q$ soit de degré total $s$ et contienne un monôme $aX_1^s$ où $\varphi(a)\ne 0$, ($\varphi$ désignant l'homomorphisme canonique $A\longrightarrow k$), que $R$ soit de degré $\leq d-s$, et que l'on ait $\bar P=\bar Q.\bar R$, (notations du chap. III, \S~4). Montrer qu'il existe alors dans $A[X_1,\ldots,X_n]$ deux polynômes $Q_{\unsure{o|0}}$, $R_o$ de degrés respectifs $s$, $d-s$, tels que $P=Q_oR_o$, $\bar Q=\bar Q_o$, $\bar R=\bar R_o$ et que $Q_o$ contienne un monôme $a_oX_1^s$ avec $\varphi(a)=\varphi(a_o)$.

(Considérer $P$, $Q$, $R$ comme des polynômes en $X_1$ à coefficients dans l'anneau $B$, complété de $A[X_2,\ldots,X_n]$ pour la valuation obtenue en prolongeant $v$ suivant la méthode du chap. VI, \S~10, n$^{\circ}$\qquad ; appliquer ensuite le lemme de Hensel ; utiliser enfin l'hypothèse initiale sur $P$).

16) Soient $B$ un anneau de valuation discrète, dont le corps résiduel $k$ est fini et n'est pas un corps premier ; soient $k_o$ le sous-corps premier de $k$, et soit $A$ le sous-anneau de $B$ formé des éléments dont la classe dans le corps résiduel appartient à $k_o$. Soient $\pi$ une uniformisante de $B$, et $(\vartheta_i)_{1\leq i\leq m}$ un système d'éléments inversibles de $B$ tels que les classes $\bar\vartheta_i\bmod.\pi$ des $\vartheta_i$
% archive sheet number: 37
\scanpage{39}
\margin{$/^{\circ}$ NB. S'il s'agit de \unsure{grpes|groupes} de type \unsure{galoisien|quelconque|galoisiens} et $U$ est ouvert \unsure{en|sur|dans} $I$, alors il s'ensuit qu'il contient un ss-groupe ouvert \unsure{plus|fin|de} $U_{\unsure{1|0}}$, qui est \uline{invariant} sous $D$}
\uline{\unsure{Sorite|Suite|Boîte} de la \unsure{moyenne|remarque}}

$D$ groupe, $I$ \unsure{ss-groupe|sous-groupe} distingué, $E$ vectoriel de dim finie sur corps $k$, sur lequel $D$ opère. On suppose que $I$ admet un sous-groupe d'indice fini $U$, \inserted{invariant \uline{dans} $D/^{\circ}$,} tel que $g\in U\Rightarrow$ $g$ opère sur $E$ de façon \uline{unipotente}, i.e. valeurs propres \unsure{ttes|toutes} $1$. Alors le groupe algébrique d'autom de $E$ engendré par les $g\in U$ est ``unipotent'', donc \deleted{\unsure{si} \unsure{on est|on a}} $E$ admet une \unsure{suite} de composition sur $U$ tel que \deleted{les} \inserted{$E$ \unsure{opère|quot.} trivialement sur les} facteurs. \deleted{\unsure{soient|sont|soit} stables sous $E$}. De façon précise, on peut prendre la suite de composition par les $E^U$, $(E/E^U)^U$ .... . Cette suite est stable par les opérations de $D$, \uline{si} \deleted{(comme il est loisible)} \deleted{on prend} \inserted{puisque} $U$ dans $I$ invariant \unsure{pour|par} les $D$. Donc il existe une suite de composition de $E$ \uline{sur $D$} tel que sur chaque facteur $E_i$, $\deleted{\unsure{I|E}}\,U$
% archive sheet number: 38
\scanpage{40}\mspage{6 - 26 -}
\hfill n$^{\circ}$ 392

à l'intérieur de chaque intervalle, prendre $b_{n+1}=b_nz^{i-1}x^{-1}$ ; à l'extrémité de chaque intervalle, utiliser l'inégalité $ijk-ij-jk-ki\geq 0$).

b) On suppose que $z^{i-1}\notin Ax+Ay$. Montrer qu'il n'existe dans l'anneau $A[[T]]$ aucune série formelle de terme constant $y^k$, ($k$ entier $>0$), et qui soit un élément associé à $v$ dans $B[[T]]$, (calculer le coefficient de $T$ dans le produit de $v$ et d'un élément inversible de $B[[T]]$).

c) Déduire de a), b) et de l'exerc. 7 un exemple d'anneau factoriel $A$ tel que l'anneau de séries formelles $A[[T]]$ ne soit pas factoriel. (Avec les notations précédentes, et en supposant remplies les conditions de a) et b) sur $x$, $y$, $z$, $i$, $j$, $k$, montrer que $vv'$ ne peut être produit de facteurs extrémaux $u_h$, $(1\leq h\leq r)$, dans $A[[T]]$ ; on observera que les $u_h$ sont des séries formelles dont les termes constants sont des puissances de $y$. Considérer ensuite l'anneau $C=S'^{-1}A[[T]]$, où $S'$ est formé des séries dont le terme constant est dans $S$ ; $B[[T]]$ est le complété de l'anneau de Zariski $C$ ; montrer que $v'\in C$, et que $v$ et les $u_h$ sont extrémaux dans $C$ ; obtenir finalement une contradiction avec b)).

d) Déduire de c) un exemple d'anneau local factoriel noethérien $A$ dont le complété $\hat A$ est un anneau de Krull non factoriel.

9) \qquad a) Soit $A$ un anneau intègre noethérien tel que pour tout idéal maximal $\underline{m}$ de $A$, $A_{\underline{m}}$ soit un anneau de valuation discrète. Si $B$ est l'anneau de séries formelles $A[[X_1,\ldots,X_n]]$, montrer que pour tout idéal maximal $\underline{n}$ de $B$, l'anneau $B_{\underline{n}}$ est factoriel, (considérer son complété, en utilisant la prop. 8 du n$^{\circ}$ 9). En déduire que tout idéal divisoriel de $B$ est un $B$-module projectif.

b) Soit $C$ un anneau noethérien tel que tout $C$-module projectif de type fini soit libre ; montrer que l'anneau de séries formelles $C[[X]]$ a la même propriété (cf. chap. II, \S~3, n$^{\circ}$ 2, prop. 5).

c) Déduire de a) et b) que si $A$ est un anneau principal, l'anneau de séries formelles $A[[X_1,\ldots,X_n]]$ est factoriel.

10) \qquad a) Un anneau factoriel prüférien, (\S~2, exerc. 12), est principal.

b) Un anneau de Krull pseudo-bezoutien, (\S~1, exerc. 21), est factoriel.

11) Soient $K$ un corps, $A$ l'anneau de polynômes $K[X,Y]$, qui est factoriel ; si $L$ est le corps $K(X^2,Y/X)\subset K(X,Y)$, montrer que l'anneau $A\cap L$ n'est pas factoriel.
% archive sheet number: 39
\scanpage{41}
opère \unsure{trivialmt|trivialement}, \inserted{(et \unsure{réciproqt|réciproquemt|respectant} bien sûr)} Donc du pt de vue virtuel, \inserted{et pour $U$ fixé,} les modules $E$ envisagés sont essentiellement les $\unsure{={}^{\mathrm{s}}|=^{\mathrm{es}}}$ que les modules tels qu'\deleted{il existe un ss-gpe \unsure{tel que|$U'$ tel que}} $U$ opère trivialmt sur $E$.

Or pour un $E$ sur lequel $U$ opère trivialmt, on peut prendre l'opération moyenne
\[
 E\leadsto E^I \quad \text{de}\quad \operatorname{Mod}(D/U,k)\ \unsure{\mathord{-}|\to}\ \operatorname{Mod}(D/I,k)
\]
et si $k$ est de car. $0$, c'est là \unsure{une|un|seul} \deleted{opération \uline{additive}, qui nous \unsure{induit|induis}} \inserted{foncteur exact. Il induit un hom} \deleted{\unsure{groupe|gr}} virtuel \deleted{$\unsure{E^I|E}$} $E\leadsto E^{\mathord{\text{\tikz[x=.66em,y=.68em,baseline=0pt]{\draw[line join=round] (.80,.95)--(.12,.24)--(.83,.34)--(.14,-.37);}}}}$
\[
 \mathord{\text{\tikz[x=.66em,y=.68em,baseline=0pt]{\draw[line join=round] (.80,.95)--(.12,.24)--(.83,.34)--(.14,-.37);}}}_U:R(D/U,k)\longrightarrow R(D/I,k)
\]

\deleted{\unsure{Ob|D}} \unsure{Obs.|On.|On a.} compatibilité, pour $V\subset U$, pour les \uline{foncteurs} $E\leadsto E^I$, d'où pour les $\mathord{\text{\tikz[x=.66em,y=.68em,baseline=0pt]{\draw[line join=round] (.80,.95)--(.12,.24)--(.83,.34)--(.14,-.37);}}}_U$ et $\mathord{\text{\tikz[x=.66em,y=.68em,baseline=0pt]{\draw[line join=round] (.80,.95)--(.12,.24)--(.83,.34)--(.14,-.37);}}}_V$ :
\diagram{\begin{tikzcd}[ampersand replacement=\&]
R(D/U,k) \arrow[r,"{\mathord{\text{\tikz[x=.66em,y=.68em,baseline=0pt]{\draw[line join=round] (.80,.95)--(.12,.24)--(.83,.34)--(.14,-.37);}}}_U}"] \arrow[d,swap,"{\text{\uline{restr.}}\,\deleted{\unsure{\mathord{\text{\tikz[x=.66em,y=.68em,baseline=0pt]{\draw[line join=round] (.80,.95)--(.12,.24)--(.83,.34)--(.14,-.37);}}}|\mathord{-}}}}"] \& R(D/I,k)\\
R(D/V,k) \arrow[ur,swap,"{\mathord{\text{\tikz[x=.66em,y=.68em,baseline=0pt]{\draw[line join=round] (.80,.95)--(.12,.24)--(.83,.34)--(.14,-.37);}}}_V}"] \&
\end{tikzcd}}
% archive sheet number: 40
\scanpage{42}\mspage{6 - 25 -}
\hfill n$^{\circ}$ 392

c) Soient $k$ un corps ordonné, $a$, $b$ des éléments $>0$ de $k$ et $B$ l'anneau $k[X,Y]/(aX^2+bY^2+1)$. Montrer que $B$ est un anneau factoriel, (utiliser b) et l'exerc. 17 a) du \S~1).

d) Montrer que l'anneau $C=\underline{Q}[X,Y]/(X^2+2Y^2+1)$ est factoriel, mais que l'anneau $C\otimes_{\underline{Q}}\underline{Q}(i)$, (avec $i^2=-1$), n'est pas factoriel.

7) Soient $K$ un anneau factoriel noethérien, $F$ un élément extrémal de l'anneau de polynômes $K[X_1,\ldots,X_n]$ ; on suppose que lorsque l'on attribue à chaque $X_i$ un poids $q(i)>0$, $(1\leq i\leq n)$, $F$ est isobare et de poids $q>0$. Soit $A$ l'anneau engendré, dans une clôture algébrique $\Omega$ de $K(X_1,\ldots,X_n)$ par $K$, les $X_i$, $(1\leq i\leq n)$, et une racine $z$ du polynôme $Z^c-F$, où $c$ est un entier étranger à $q$. Montrer que $A$ est factoriel dans les deux cas suivants :

1) on a $c\equiv 1\ (\mathrm{mod.}\ q)$ ;

2) tout $K$-module projectif de type fini est libre, (ce qui a lieu par exemple lorsque $K$ est un corps, ou un anneau principal, ou un anneau local).

(Dans le premier cas, considérer l'anneau de fractions $A[1/z]$ ; montrer qu'il est factoriel et appliquer la prop. 3 du n$^{\circ}$ 4. Dans le second cas, considérer un entier $d$ tel que $cd\equiv 1\ (\mathrm{mod.}\ q)$, et soit $z'\in\Omega$ tel que $z=z'^{d}$ ; l'anneau $B=A[z']$ est factoriel en vertu du premier cas, et est un $A$-module libre. Considérer $B$ comme un anneau gradué en prenant $z'$ de poids $q$, chacun des $X_i$ de poids $cdq(i)$, et se ramener à prouver que pour deux éléments homogènes $u$, $v$ de $A$, l'idéal $Au\cap Av$ est principal, en utilisant l'exerc. 16 du \S~1 ; considérer enfin l'idéal $Bu\cap Bv$ dans $B$ et utiliser le chap. I, \S~3, n$^{\circ}$ 6, prop. 12). En particulier, si $K$ est un corps, $a$, $b$, $c$ trois entiers $>0$ étrangers deux à deux, l'anneau $A=K[X,Y,Z]/(Z^a-X^b-Y^c)$ est factoriel.

8) Soient $A$ un anneau intègre, $x$, $y$, $z$ trois éléments non nuls de $A$ où $x$ est extrémal et $Ax\cap Ay=Axy$. Soit $S$ l'ensemble multiplicatif des $x^n$, $(n\geq 0)$, et soit $B=S^{-1}A$ ; on considère les anneaux de séries formelles $A[[T]]$ et $B[[T]]$.

a) Soient $i$, $j$, $k$ trois entiers $\geq 0$ tels que $ijk-ij-jk-ki\geq 0$ et $z^i\in Ax^j+Ay^k$. On considère dans $A[[T]]$ l'élément $v=xy-z^{i-1}T$. Montrer qu'il existe un entier $t>0$ et une série $v'=y^t x^{-1}+b_1x^{-2}T+\ldots$ $\ldots+b_{n-1}x^{-n}T^{n-1}+\ldots$ dans $B[[T]]$ tels que $vv'\in A[[T]]$, (on déterminera par récurrence les $b_n$, en procédant dans $\underline{N}$ par intervalles de longueur $ij$ ;
% archive sheet number: 41
\scanpage{43}
D'autre part, on a
\[
 R(D/U,k)\hookrightarrow R(D,k)
\]
et l'image de $R(D/U,k)$ dans $R(D,k)$ est précisément \unsure{formé|formée} des classes des $E$ \unsure{satisfaisant|satisfaisants} la condition du début. Pour un tel $E$, on a donc une opération $\mathord{\text{\tikz[x=.66em,y=.68em,baseline=0pt]{\draw[line join=round] (.80,.95)--(.12,.24)--(.83,.34)--(.14,-.37);}}}$ sur la \uline{classe} de $E$. Donc $\operatorname{cl}(E)^{\mathord{\text{\tikz[x=.66em,y=.68em,baseline=0pt]{\draw[line join=round] (.80,.95)--(.12,.24)--(.83,.34)--(.14,-.37);}}}}$ est défini dès qu'il existe un ss-gpe \deleted{\unsure{d'ind|d'indice}} $U$ d'indice fini dans $I$ tel que $\deleted{D/U}$ $U$ opère sur $E$ de façon unipotente. / Si on prend des classes de repr. \uline{continues} ($k$ étant un corps top.) on \deleted{\unsure{parler|parlera|prendra}} \inserted{\unsure{obtient|obtiendra}} de $=$ $\operatorname{cl}(E)^{\mathord{\text{\tikz[x=.66em,y=.68em,baseline=0pt]{\draw[line join=round] (.80,.95)--(.12,.24)--(.83,.34)--(.14,-.37);}}}}$ comme classe de repr. \uline{continues} de $D/I$. Compatibilité avec ext du corps de base.

\deleted{Supposons que $D/I\simeq\mathbb{Z}$, \unsure{et notons|et disons|et soit|et donc toute} qu'une classe de repr. sur $D/I$ est \unsure{donné|donnée|déterminé|déterminée} par \unsure{sa} série car. $-$ $L_{E^{\mathord{\text{\tikz[x=.66em,y=.68em,baseline=0pt]{\draw[line join=round] (.80,.95)--(.12,.24)--(.83,.34)--(.14,-.37);}}}}}(t)$. \unsure{On|Ceci|Comme} \unsure{aura|arrive} alors}
% archive sheet number: 42
\scanpage{44}\mspage{6 - 24 -}
\hfill n$^{\circ}$ 392

4) \qquad a) Montrer que l'anneau $A=\underline{Q}[X,Y]/\underline{p}$, où $\underline{p}$ est l'idéal principal engendré par $X^2+Y^2-1$ dans $\underline{Q}[X,Y]$, est un anneau de Krull non factoriel, (montrer que $X$ est un élément extrémal de $A$ mais que $AX$ n'est pas un idéal maximal de $A$).

b) Montrer que l'anneau $A\otimes_{\underline{Q}}\underline{Q}(i)$ est factoriel, (prouver que cet anneau est isomorphe au quotient de $\underline{Q}(i)[X,Y]$ par l'idéal principal $(XY-1)$ ; comparer à l'exerc. 3 b)).

5) \qquad a) Soient $k$ un anneau factoriel noethérien, $B$ l'anneau de polynômes $k[X_1,\ldots,X_n]$ avec $n\geq 3$ ; soient $g_i$, $(0\leq i\leq r)$, des éléments de $k[X_3,\ldots,X_n]$ où $g_{\unsure{o|0}}$ est extrémal. On pose $g=X_1X_2-\sum_{i=o}^{r}g_iX_1^{i}$, et on considère l'anneau quotient $A=B/gB$. Montrer que $A$ est factoriel. (Soit $S$ la partie multiplicative de $A$ engendrée par $1$ et l'image de $X_1$ dans $A$ ; appliquer à $S^{-1}A$ la prop. 3 du n$^{\circ}$ 4).

b) Soient $k$ un corps de caractéristique $\ne 2$, $F$ un polynôme homogène du second degré de $k[X_1,\ldots,X_n]$ avec $n\geq 5$, tel que la fonction polynôme correspondante sur $k^n$ soit une forme quadratique non dégénérée. Montrer que l'anneau $k[X_1,\ldots,X_n]/(F)$ est factoriel. (On prouvera d'abord qu'un polynôme homogène $G$ du second degré dans $k[X_1,\ldots,X_n]$ tel que la fonction polynôme correspondante soit une forme quadratique non dégénérée, est extrémal pour $n\geq 3$. Prouver ensuite la proposition lorsque $k$ est algébriquement clos, en utilisant a) ; passer enfin au cas général à l'aide de l'exerc. 3 b)).

c) Si $F=X_1X_2-X_3X_4$, montrer que l'anneau $k[X_1,X_2,X_3,X_4]/(F)$ n'est pas factoriel. (Montrer que les images des $X_i$ dans cet anneau sont des éléments extrémaux).

6) \qquad a) Soit $K$ un corps algébriquement clos de caractéristique $\ne 2$. Déterminer les idéaux premiers gradués de l'anneau $K[X,Y,Z]/(X^2+Y^2+Z^2)$.

b) Soient $k$ un corps ordonné, $a$, $b$, $c$ des éléments $>0$ de $k$ et $A$ l'anneau $k[X,Y,Z]/(aX^2+bY^2+cZ^2)$. Montrer que $A$ est un anneau factoriel. (Se ramener au cas où $k$ est un corps ordonné maximal, en utilisant l'exerc. 3 b) ; prouver alors, en utilisant a), que tout idéal premier gradué de hauteur 1 dans $A$ est principal, et appliquer l'exerc. 3 a)).
% archive sheet number: 43
\scanpage{45}
Calculons la \uline{trace} de la représentation \unsure{dévissée|divisée|ci-dessus}, on trouve pour $\varphi\deleted{\unsure{\xi|g}}\in D/I$
\[
 \operatorname{Tr}_{E^{\mathord{\text{\tikz[x=.66em,y=.68em,baseline=0pt]{\draw[line join=round] (.80,.95)--(.12,.24)--(.83,.34)--(.14,-.37);}}}}}(\varphi)\deleted{\unsure{\mathord{\text{\tikz[x=.66em,y=.68em,baseline=0pt]{\draw[line join=round] (.80,.95)--(.12,.24)--(.83,.34)--(.14,-.37);}}}|g}}
 =\sum_i\operatorname{Tr}_{E_i^{\mathord{\text{\tikz[x=.66em,y=.68em,baseline=0pt]{\draw[line join=round] (.80,.95)--(.12,.24)--(.83,.34)--(.14,-.37);}}}}}(\varphi)
\]

Or si $\deleted{\unsure{I|G}}$ $U$ opère trivialmt sur $E$, \deleted{qui est} alors un module sur $D/U$, on a
\[
 \deleted{\unsure{\text{en}|\text{on}}}\ E^{\mathord{\text{\tikz[x=.66em,y=.68em,baseline=0pt]{\draw[line join=round] (.80,.95)--(.12,.24)--(.83,.34)--(.14,-.37);}}}}
 =\operatorname{Im}\deleted{\unsure{\tfrac{1}{N}|\mathord{\text{\tikz[x=.66em,y=.68em,baseline=0pt]{\draw[line join=round] (.80,.95)--(.12,.24)--(.83,.34)--(.14,-.37);}}}}}\underset{\text{\inserted{\unsure{projection|projecteur|régulière}}}}{\pi}\deleted{\unsure{E|I}}
 \qquad \pi=\frac{1}{N}\sum_{i\in I/U}e_i \quad .\ N=\operatorname{Card} I/U
\]
donc si $g\in I/U\longrightarrow\varphi$ \unsure{on aura|on a|\ldots\ \ldots}
\[
 \operatorname{Tr}_{E^{\mathord{\text{\tikz[x=.66em,y=.68em,baseline=0pt]{\draw[line join=round] (.80,.95)--(.12,.24)--(.83,.34)--(.14,-.37);}}}}}(\varphi)=\operatorname{Tr}_E(\deleted{\unsure{I|1}}\pi g)
 \deleted{=\frac{1}{N}\sum_{\substack{g\to\varphi\\\unsure{\text{dans }I|g\in I}}}\operatorname{Tr}_E(}
 \qquad\text{d'où}
\]
\[
 \boxed{\displaystyle\operatorname{Tr}_{E^{\mathord{\text{\tikz[x=.66em,y=.68em,baseline=0pt]{\draw[line join=round] (.80,.95)--(.12,.24)--(.83,.34)--(.14,-.37);}}}}}(\varphi)=
 \frac{1}{N}\sum_{\substack{\psi_\alpha\in D/U\\\unsure{\substack{\deleted{i}\\\psi_\alpha}|\bar{\psi}_\alpha|\substack{\mathord{-}\\\psi_\alpha}|\psi_\alpha}\to\varphi\in D/I}}\operatorname{Tr}_E\psi_\alpha}
\]

Comme \inserted{les deux \unsure{mbres|membres} de} cette formule sont linéaires en $E$, ça reste valable pour tt $E$ ; il faut simplemt noter que $\operatorname{Tr}_E\psi$ est défini \deleted{\unsure{même si|si}} pour $\psi\deleted{\unsure{\alpha|i}}\in D/U$, si $U$ opère de façon unipotente, comme $\operatorname{Tr}_E(\unsure{\tau|\gamma|\psi}')$ pour $\tau'\in D$ se restreignant en $\tau$ : cette valeur ne \uline{dépend pas} du
% archive sheet number: 44
\scanpage{46}\mspage{4 - 42 -}
\hfill n$^{\circ}$ 392

morphe à un idéal de $A$. Supposons donc $n\geq 2$ ; en vertu du lemme 8, il existe un sous-module libre $L_{\unsure{o|0}}$ de $M$, de rang 1, tel que $M/L_o$ soit sans torsion ; $M/L_o$ est donc de rang $n-1$. En vertu de l'hypothèse de récurrence, il y a donc un sous-module libre $L_1$ de $M/L_o$ tel que $(M/L_o)/L_1$ soit isomorphe à un idéal de $A$. Soit $L$ l'image réciproque de $L_1$ dans $M$ ; $L/L_o$ est isomorphe à $L_1$, et puisque $L_1$ est libre, $L$ est isomorphe à $L_o\oplus L_1$, (\uline{Alg}., chap. II, 3e éd., \S~1, n$^{\circ}$ 11, prop. 21), donc libre ; comme $M/L$ est isomorphe à $(M/L_o)/L_1$, le théorème est démontré.

\uline{Remarque}.- Si $M$ est réflexif, il n'en est pas nécessairement de même de $M/L$, (exerc. \qquad).

10. \uline{Modules sur les anneaux de Dedekind}.

On suppose maintenant que $A$ soit un \uline{anneau de Dedekind} ; on sait alors que les idéaux $\underline{p}\in P$ sont \uline{maximaux} et que ce sont les seuls idéaux premiers $\ne 0$ de $A$, (\S~2, n$^{\circ}$ 2, th. 1) ; le groupe $D(A)$ s'identifie au groupe $I(A)$ des idéaux fractionnaires $\ne 0$ de $A$.

PROPOSITION 21.- \uline{Soit} $A$ \uline{un anneau de Dedekind}. \uline{Tout} $A$-\uline{module pseudo-nul est nul}. \uline{Tout homomorphisme pseudo-injectif}, (resp. \uline{pseudo-surjectif}, \uline{pseudo-bijectif}, \uline{pseudo-nul}), \uline{de} $A$-\uline{modules est injectif}, (resp. \uline{surjectif}, \uline{bijectif}, \uline{nul}).

La première assertion a déjà été démontrée, (n$^{\circ}$ 4, \uline{Exemple} 1) ; les autres s'en déduisent immédiatement.

PROPOSITION 22.- \uline{Soient} $A$ \uline{un anneau de Dedekind}, $M$ \uline{un} $A$-\uline{module de type fini}. \uline{Les propriétés suivantes sont équivalentes} :

a) $M$ \uline{est sans torsion}.

b) $M$ \uline{est réflexif}.

c) $M$ \uline{est projectif}.

On sait déjà, (sans hypothèse sur l'anneau intègre $A$), que b) implique a), (n$^{\circ}$ 2, \uline{Remarque} 1), et que c) implique b), (\uline{Alg}., chap. II, 3e éd., \S~2, n$^{\circ}$ 7,
% archive sheet number: 45
\scanpage{47}
choix de $\unsure{\tau|\gamma|\psi}'$ relevant $\tau$, comme on voit à l'aide de la suite de composition.

\begin{center}\rule{4em}{.4pt}\end{center}

Revenons au cas géométrique, on prendra pour $D$ le gpe de décomposition $D_x$ d'un $x\in X$, $I$ le gpe d'inertie, d'où $D_x/I_x\simeq\pi_x\simeq\widehat{\mathbb{Z}}$. (générateur $f_x$). On posera
\[
 E_{\unsure{\eta|y}}(x)=E_\eta^{\mathord{\text{\tikz[x=.66em,y=.68em,baseline=0pt]{\draw[line join=round] (.80,.95)--(.12,.24)--(.83,.34)--(.14,-.37);}}}(x)}\qquad\text{comme module sur }\pi_x.
\]

\deleted{On posera} Il revient au $=$, du pt de vue faisceau virtuel sur $X$, de poser
\[
 E_\eta^{\mathord{\text{\tikz[x=.66em,y=.68em,baseline=0pt]{\draw[line join=round] (.80,.95)--(.12,.24)--(.83,.34)--(.14,-.37);}}}}=i_!(E_U)+\sum_{x\in Y}j_{x*}\bigl(E_\eta^{\mathord{\text{\tikz[x=.66em,y=.68em,baseline=0pt]{\draw[line join=round] (.80,.95)--(.12,.24)--(.83,.34)--(.14,-.37);}}}(x)}\bigr)
\]
où $U$ est un ouvert dense de $X$ sur lequel $E_\eta$ n'est pas ramifié, $E_U$ est l'image de $E_\eta$ sur $U$, \inserted{$Y=X-U$,} $i:U\mathord{-}X$ \inserted{et $j:Y\mathord{-}X$} les immersions canoniques, $j_x:\operatorname{Spec}k(x)\longrightarrow X$.
% archive sheet number: 46
